\section{Developing the scanner and parser in F\#}
Integer-C is the input language. The compiler's host language is F\# and its target is JVM/Jasmin assembly.
The grammar and Integer-C sample above therefore retain their original notation.
The frontend turns that syntax into typed records and union cases.
@@snippet:compiler-types@@

The FsLex specification recognizes keywords, identifiers, constants, punctuation and comments.
Longer comparison operators are recognized as complete tokens. Comment scanning is a separate rule;
unclosed comments and illegal characters carry line/column diagnostics.
@@snippet:compiler-lexer@@

The FsYacc specification constructs the typed AST. Empty parameter, declaration and statement lists return actual
empty F\# lists. Function definitions retain their name, parameters, local declarations and body in separate fields.
Precedence rules and the grammar structure handle arithmetic and comparisons, and the dangling else belongs
to the nearest unmatched conditional.
@@snippet:compiler-parser@@

The frontend boundary creates the buffer, initializes its filename/position and turns scanner/parser failures into
a source location. Generated internals are not exposed as the compiler's public parse API.
@@snippet:compiler-frontend@@

\href{../../Chapter-12/Compiler.ipynb}{\nolinkurl{Chapter-12/Compiler.ipynb}}
prints the token stream and AST before code generation. This separates lexical errors, syntax errors and backend errors.
The source examples include the original sum, factorial and prime-number programs, plus precedence and ternary-expression checks.
